Budgeted teams of mobile agents can cover more spatially distributed tasks than one agent, but independent routes may prevent timely data sharing. Team orienteering captures reward selection and per-agent budgets~\cite{chao1996team}. Kinematic orienteering adds bounded velocity and acceleration to the flight-time calculation~\cite{meyer2022kop}. In missions such as disaster assessment, agents may need to separate temporarily and communicate later, so a requirement for connectivity at every instant can be restrictive.

We consider routes that start and end at common depots, collect reward only at explicitly declared task visits, and may include freely positioned visits with no reward. Each agent has a flight-time budget that includes motion and pre-arrival waiting. Trajectories are continuous in time. A pairwise link becomes usable after a continuous within-range handshake, and information can be relayed both simultaneously and through contacts at later times. The communication requirement is a bound $T_c$ on the worst directed information age during the mission. This is stronger than connectivity of a union of contact graphs over a window: a sequence of contacts in the wrong order need not deliver the information.

Our contribution is an integrated formulation and implementation of this reward--motion--communication problem. We extend the kinematic orienteering representation to a team with shared reward accounting, add continuous-time contact extraction and chronological information propagation, and search over routes, task assignments, meeting structure, and free waypoints with a multi-objective variable-neighborhood method. The three minimized objectives are negative team reward, mean completion time, and positive violation of $T_c$. Feasible solutions therefore compete on reward and time without an artificial incentive for connectivity beyond the required bound.

We report reproducible preliminary experiments from saved solver outputs. Four stylized maps test clustered, elongated, branching, and uniform task layouts with three agents. A separate two-agent probe isolates the potential value of a free waypoint. The experiments reveal both feasible plans and a pronounced search limitation: within a 115-second optimizer allowance, the solver often retains convoy-like routes and rarely reaches its communication neighborhoods on the larger maps. We do not claim a matched TOP-CC benchmark or optimality result.

Section~\ref{sec:related} places the model alongside orienteering and intermittent-connectivity planning. Sections~\ref{sec:problem_formulation} and~\ref{sec:method} state the implemented model and search method. Section~\ref{sec:experiments} reports the saved experiments and their limitations.
